\chapter{Service Architecture}
\label{chap:service-architecture}

Protection domains contain service faults and authority. Naming and protocol
policy remain in userspace; privileged code manages lifecycle and resources.

\section{Protection domains as failure boundaries}

A protection domain (\autoref{chap:vocabulary}) is the unit of both authority and failure. When
a service crashes:

\begin{itemize}
    \item Its capability table is destroyed, invalidating its handles --- client
        connections to the closed endpoint fail through the IPC lifecycle.
    \item Its memory mappings are removed --- outstanding borrows are revoked; the
        kernel unmaps the pages from the (dead) borrower and restores the lender's
        access.
    \item Its threads are aborted --- any thread blocked on a CQ or endpoint is
        woken with a terminal error.
    \item Its device grants are recovered --- the supervisor regains ownership of
        any MMIO region and interrupt objects the domain held.
\end{itemize}

The MMU and capability checks are intended to contain the fault to that domain;
kernel, hardware, and DMA defects remain outside this guarantee.

\section{The supervisor --- domain lifecycle manager}

The \textbf{supervisor} is a kernel-side component that creates and destroys
protection domains. It does not understand service protocols, but it does understand
ELF images, memory mappings, and capability grants.

Its operations are:

\begin{description}
    \item[Load] Parse an ELF image, validate its segments, create a
        new address space, map \texttt{PT\_LOAD} segments at their linked
        virtual addresses, set up a stack with a guard page.
    \item[Grant] Populate the domain's config-page with the bootstrap
        capabilities it needs:
        \begin{itemize}
            \item A connection to the name service (for service lookup).
            \item For driver domains: an \texttt{MmioRegion} capability, an
                \texttt{Interrupt} capability.
            \item For service domains: an \texttt{Endpoint} capability on which
                to listen.
            \item A default CQ handle.
        \end{itemize}
    \item[Spawn] Create the initial thread with the domain's entry
        point and its affinity LP.
    \item[Monitor] Register to be notified when the domain exits.
    \item[Teardown] on exit, reclaim all capabilities, revoke
        mappings, reconcile outstanding operations, and make the address-space
        ID available for reuse.
\end{description}

Grants are minted kernel-side by the supervisor and never through a syscall. A
driver does not ask the kernel for MMIO; the supervisor grants it at creation
time based on the device topology it enumerated during boot.

\section{The name service --- discovery without the kernel}

The kernel does not know about service names. A userspace \textbf{name service}
maps human-readable names to connection capabilities:

\begin{description}
    \item[Registration] A service calls \texttt{register("driver.nic.v1",
        endpoint, access\_key)} on the name service. The name service stores the
        endpoint's identity and opens a connection to it.
    \item[Generation tracking] Each registration bumps a generation counter.
        If the service restarts and re-registers with the same name, the new
        generation replaces the old one. Existing connections bound to the old
        generation fail with \texttt{ServiceRestarted}.
    \item[Lookup] A client calls \texttt{lookup("driver.nic.v1", access\_key)}.
        The name service checks the access key (if the service registered one),
        then delegates a connection capability to the client --- typically with
        \texttt{SEND | CALL} rights only.
    \item[Access-key gating] A registration with a non-zero access key requires
        lookers to provide the matching key. This is userspace policy: the name
        service enforces it; the kernel does not participate.
\end{description}

The name service runs as a \texttt{no\_std} EL0 service process. Boot starts
exactly one node-local instance and the supervisor retains its registry
endpoint. Each ordinary service or application receives an attenuated
connection to that endpoint through its bootstrap page, so all processes on
the node observe the same registry. Registering another name does not create
another name-service process. A separate registry is appropriate only for an
explicitly isolated namespace, such as a test or a future tenant domain.

The supervisor's endpoint handle remains kernel-private. Possessing a
bootstrap connection permits the operations granted on that connection; it
does not permit an application to mint arbitrary registry connections or
inspect another process's capability table.

\section{Service restart and recovery}

When a service crashes and the supervisor restarts it:

\begin{enumerate}
    \item Device reset --- if the service was a driver, the supervisor
        resets the device before re-granting it to a new driver instance. The
        new driver sees a clean hardware state.
    \item Generation increment --- the name service bumps the service's
        generation. All old connections are invalidated.
    \item Stale connections --- clients connected to the old instance
        receive \texttt{ServiceRestarted} on their next call
        attempt. They re-lookup the service and receive a connection to the
        new instance.
    \item Outstanding operation reconciliation --- the supervisor
        reconciles any operations the old domain had submitted but not
        completed: cancellations are posted, borrowed memory is returned,
        DMA buffers are recovered.
    \item Fresh bootstrap --- the new domain receives a fresh config-page
        with freshly minted capabilities. It does not inherit the previous
        instance's capability table, reply tokens, or DMA ownership.
\end{enumerate}

A restarted service receives fresh authority. Tests cover stale connections and
operation reconciliation; the status appendix records the verified scope.

\section{The overall system topology}

At steady state, a CharlotteOS system consists of:

\begin{itemize}
    \item The kernel --- capabilities, endpoints, memory objects,
        CQs, scheduling, interrupt routing.
    \item The supervisor (kernel-side) --- domain lifecycle, ELF
        loading, device enumeration.
    \item  The node name service (one userspace instance) --- the
        shared registry for service registration, lookup, access-key gating,
        and generation tracking.
    \item Service processes --- one or more per logical function:
        NIC driver, TCP/IP stack, UTC time, filesystem, Raft consensus,
        application servers.
    \item Client processes --- application code that looks up
        services and invokes them through capabilities.
\end{itemize}

Every component beyond the kernel, supervisor, and name service is a replaceable
userspace process. A different NIC driver, a different filesystem, a different
network stack --- these are configuration choices, not kernel modifications.

\section{Why service architecture matters for critical software}

\begin{enumerate}
    \item Failure containment is architectural --- a service fault closes
        its endpoints and affects clients through defined errors rather than
        direct memory corruption. Availability still depends on the failed
        service and its dependants.

    \item Recovery follows a defined sequence --- the restart sequence
        (device reset, generation bump, connection invalidation, operation
        reconciliation) is the same for every driver. A new driver
        implemented to the same endpoint protocol replaces the old one
        without any special-case recovery code.

    \item Policy lives in userspace --- the kernel enforces mechanism
        (capabilities, endpoints, memory objects). Userspace enforces policy
        (who can register what name, who can look up what service, how
        many restarts are allowed, what log level to use). The kernel
        stays small; policy is auditable and replaceable.

    \item The bootstrap contract is explicit and minimal --- every
        domain starts with exactly the capabilities it needs and no more.
        The config-page contract defines what a domain receives; there is no
        ambient authority to discover or exploit.
\end{enumerate}

\endinput
